\section{Related Works: Autonomous Endoscopy and Robotic Polypectomy}
\label{sec:related_works}

In recent years, the intersection of \textbf{computer vision}, \textbf{reinforcement learning}, and \textbf{robotic endoscopic procedures} has seen rapidly growing interest. Early research mostly looked into supervised deep learning for passive, real-time polyp detection. These models got great sensitivity, but they couldn't actually do anything physical in the real world \cite{wang2019cade,repici2020cade,nazarian2021jmir}. Now, however, the field is moving towards active and autonomous intervention.

\textbf{Deep Reinforcement Learning (DRL) for Navigation:}
Several projects have attempted to automate or assist in colonoscopy navigation. For instance, projects like those by Pore et al. demonstrated that \textbf{DRL agents} could learn basic navigation and end-to-end visuomotor control in simplified tubular environments using purely visual inputs \cite{pore2022iros}. Similarly, Corsi et al. explored \textbf{constrained reinforcement learning} combined with formal verification to guarantee safe colonoscopy navigation, ensuring the endoscope avoids dangerous collisions with the mucosal walls \cite{corsi2023arxiv}. 

\textbf{Surgical Robotics and Simulation:}
Beyond pure navigation, the development of sophisticated simulation environments has been a catalyst for surgical robot learning. The \textbf{SurRoL} platform, for example, provides an open-source, reinforcement learning-centered simulation environment compatible with the dVRK (da Vinci Research Kit), facilitating the training of agents for complex tissue manipulation \cite{surrol2021arxiv}. Researchers such as Ou and Tavakoli have leveraged similar sim-to-real frameworks to train autonomous planning agents capable of internal tissue manipulation using RL \cite{ou2023lra}. Additionally, comprehensive reviews like those by Qian et al. emphasize how DRL is enhancing the automation level in surgical robotics, moving systems from mere teleoperation toward semi-autonomous execution \cite{qian2023arxiv}.

\textbf{Our Contribution:}
However, despite these advances, existing works often struggle with the complex, deformable physics of the colon and the highly intricate, multi-step process of \textbf{polyp resection}. Navigating a rigid tube is vastly different from safely deploying a snare, capturing a polyp, and executing a cut without causing thermal injury or perforation. Our work directly builds upon these foundations. By integrating a \textbf{CNN visual encoder} with an \textbf{LSTM-equipped PPO agent}, we tackle not just safe spatial navigation, but the delicate, sequential task of \textbf{snare capture and polypectomy} inside a highly realistic, dynamic virtual environment. We aim to bridge the gap between simple lumen following and complete, autonomous lesion extraction.
